\documentclass[a4paper]{article}

% Basic packages
% Español (México) con XeLaTeX: fontspec para fuentes Unicode, polyglossia
% para la división silábica y los nombres del documento en la variante mexicana.
\usepackage{fontspec}
\usepackage{polyglossia}
\setdefaultlanguage[variant=mexican]{spanish}
% Los nombres chinos van como «pinyin (original)»: xeCJK da a los caracteres
% Han su propia fuente; sin ella XeLaTeX los omite con un aviso.
% FandolSong no cubre algunos caracteres de nombres (U+73FA); WenQuanYi Zen Hei sí.
\usepackage[AutoFallBack=true]{xeCJK}
\setCJKmainfont{FandolSong-Regular.otf}
\setCJKfallbackfamilyfont{\CJKrmdefault}{WenQuanYi Zen Hei}
% polyglossia no traduce los nombres de listings.
\AtBeginDocument{\providecommand{\lstlistingname}{}\renewcommand{\lstlistingname}{Listado}%
  \providecommand{\lstlistlistingname}{}\renewcommand{\lstlistlistingname}{Índice de listados}}
\usepackage{amsmath, amssymb}
\usepackage{graphicx}
\usepackage[margin=2.5cm]{geometry}
\usepackage[most]{tcolorbox}
\usepackage{etoolbox}
\usepackage{listings}
\usepackage{hyperref}
\usepackage{booktabs}
\usepackage{subcaption}
\usepackage{float}

% Highlight boxes
\newtcolorbox{knowledgebox}[1]{
    enhanced, colback=blue!5!white, colframe=blue!75!black, colbacktitle=blue!75!black,
    coltitle=white, fonttitle=\bfseries, title={#1}, attach boxed title to top left={yshift=-2mm, xshift=2mm},
    boxrule=1pt, sharp corners
}
\newtcolorbox{importantbox}[1]{
    enhanced, colback=yellow!10!white, colframe=yellow!80!black, colbacktitle=yellow!80!black,
    coltitle=black, fonttitle=\bfseries, title={#1}, sharp corners
}
\newtcolorbox{warningbox}[1]{
    enhanced, colback=red!5!white, colframe=red!75!black, colbacktitle=red!75!black,
    coltitle=white, fonttitle=\bfseries, title={#1}, sharp corners
}

\lstset{
    language=Python, basicstyle=\ttfamily\small, keywordstyle=\color{blue},
    stringstyle=\color{red!60!black}, commentstyle=\color{green!60!black},
    breaklines=true, frame=single, numbers=left, numberstyle=\tiny\color{gray}, captionpos=b, extendedchars=true
}

% Metadata
\newcommand{\notetitle}{Notas de la entrevista: Dario Amodei, Amanda Askell \& Chris Olah}
\newcommand{\noteauthors}{Elaboradas a partir de materiales públicos del curso}
\newcommand{\notedate}{2025-03-26}
\newcommand{\videochannel}{Lex Fridman Podcast}
\newcommand{\videopublishdate}{2025-03-26}
\newcommand{\videoduration}{5 horas 15 minutos}
\newcommand{\videourl}{https://www.youtube.com/watch?v=hFZFjoKJR0I}
\newcommand{\videocoverpath}{cover.jpg}
\newcommand{\repourl}{https://github.com/hqhq1025/ai-course-notes}

\usepackage{fancyhdr}
\pagestyle{fancy}
\fancyhf{}
\fancyfoot[C]{\thepage}
\fancyfoot[R]{\footnotesize\color{gray}\href{\repourl}{AI Course Notes}}
\renewcommand{\headrulewidth}{0pt}
\renewcommand{\footrulewidth}{0pt}

\begin{document}

\begin{titlepage}
\centering
{\Large Notas de la entrevista\par}
\vspace{1.2cm}
{\huge\bfseries \notetitle\par}
\vspace{0.8cm}
{\large \noteauthors\par}
\vspace{0.3cm}
{\large \notedate\par}
\vspace{1.2cm}
\ifdefempty{\videocoverpath}{\vspace{4cm}}{
\includegraphics[width=0.82\textwidth,height=0.45\textheight,keepaspectratio]{\videocoverpath}\par
}
\vfill
\begin{tcolorbox}[width=0.9\textwidth, colback=black!2!white, colframe=black!60, sharp corners]
\textbf{Canal}: \videochannel\par
\textbf{Fecha de publicación}: \videopublishdate\par
\textbf{Duración del video}: \videoduration\par
\ifdefempty{\videourl}{}{\textbf{Enlace del video}: \href{\videourl}{\nolinkurl{\videourl}}\par}\par
\textbf{Página del proyecto}: \href{\repourl}{\nolinkurl{\repourl}}\par
\end{tcolorbox}
\vspace{0.3cm}
{\small Esta entrevista incluye a tres invitados: Dario Amodei (director ejecutivo de Anthropic, aproximadamente 2.5 horas), Amanda Askell (diseño de la personalidad de Claude, aproximadamente 1.5 horas) y Chris Olah (interpretabilidad mecanicista, aproximadamente 1 hora).\par}
\end{titlepage}

\tableofcontents
\newpage

%% ==================== Part I: Dario Amodei ====================
\part{Dario Amodei: Scaling, Claude y AGI}

\section{Scaling Laws: de la intuición empírica a una convicción central}

\subsection{Origen: la revelación de 2014-2017}

Dario se incorporó a Baidu en 2014, bajo la mentoría de Andrew Ng, para trabajar en reconocimiento de voz. En ese entonces, la opinión dominante en el mundo académico era que «nos falta el algoritmo correcto», pero Dario, como recién llegado, planteó una pregunta sencilla: ¿qué pasaría si el modelo se hiciera más grande, se agregaran más datos y se entrenara durante más tiempo?

\begin{knowledgebox}{El origen de la Scaling Hypothesis}
Dario no es el único creyente en el Scaling. Ilya Sutskever tiene una frase célebre y casi zen: \textit{«The thing you need to understand about these models is they just want to learn.»}

El Bitter Lesson de Rich Sutton y la Scaling Hypothesis de Gwern apuntaban en la misma dirección. Pero el verdadero punto de inflexión fueron los resultados de GPT-1 en 2017, que convencieron a Dario de que el lenguaje era el mejor vehículo para el Scaling.
\end{knowledgebox}

\begin{importantbox}{Analogía de la reacción química}
El Scaling requiere ampliar linealmente y al mismo tiempo tres «reactivos»: una red más grande, más datos y más cómputo. \textit{«Si solo amplías uno de ellos sin ampliar los demás, la reacción se detiene.»} Esto explica por qué muchos investigadores, en las primeras etapas, no vieron el efecto del Scaling: solo ampliaron una o dos de esas dimensiones.
\end{importantbox}

\subsection{¿Por qué funciona el Scaling? El ruido 1/f y la distribución de cola larga}

Dario recurre a su formación en física (en el posgrado estudió biofísica) para explicar el mecanismo subyacente del Scaling. Introduce un concepto clave: el \textbf{ruido 1/f}, es decir, la idea de que los procesos naturales producen distribuciones de cola larga.

El lenguaje natural tiene una estructura jerárquica que va de lo simple a lo complejo:
\begin{enumerate}
    \item Distribución de frecuencia de palabras (los patrones más comunes)
    \item Reglas gramaticales básicas
    \item Estructura a nivel de oración
    \item Coherencia temática a nivel de párrafo
    \item Ideas novedosas y razonamiento complejo
\end{enumerate}

\textit{«Primero capturan correlaciones muy simples… y después viene una cola muy larga con otros patrones más complejos.»} La suavidad de la distribución de cola larga se refleja directamente en el descenso suave de la curva de loss: por eso las Scaling Laws muestran una relación de ley de potencia tan regular.

Idea clave: este marco no solo explica \textit{por qué} funciona el Scaling, sino también \textit{por qué el progreso parece gradual y no repentino}: porque el modelo va «excavando» progresivamente a lo largo de la distribución de cola larga hacia patrones cada vez más raros.

\subsection{¿Dónde está el techo de la inteligencia?}

Dario distingue tres niveles de techo:

\textbf{Por debajo del nivel humano}: no existe techo. \textit{«Los humanos somos capaces de entender estos patrones»}, así que en teoría el modelo también debería poder alcanzarlos.

\textbf{Más allá del nivel humano: depende del campo}. En biología puede haber un margen enorme: \textit{«los humanos apenas logramos entender la complejidad de la biología»}, un campo con disciplinas muy fragmentadas donde es difícil integrar una visión de conjunto. En cambio, en campos como la ciencia de materiales o el conflicto humano, el techo podría estar más cerca del nivel humano.

\textbf{El techo institucional}: esta es la opinión más singular de Dario:

\begin{warningbox}{El verdadero cuello de botella no es la inteligencia, sino las instituciones}
El sistema de ensayos clínicos es un ejemplo excelente: incluye tanto burocracia innecesaria (procesos excesivamente conservadores) como mecanismos de protección razonables (la seguridad de los pacientes). Aunque la IA lograra diseñar un fármaco perfecto, este igual tendría que pasar por ese proceso lento.

\textit{«Somos demasiado lentos y demasiado conservadores.»} Dario considera que las instituciones humanas (los sistemas burocráticos, los sistemas regulatorios, la inercia organizacional) podrían ser un cuello de botella mayor que las limitaciones puramente de inteligencia.
\end{warningbox}

\subsection{Posibles obstáculos para el Scaling}

Dario analiza uno por uno los factores que podrían impedir que el Scaling continúe:

\textbf{Limitaciones de datos}: los datos disponibles en internet son limitados y su calidad está bajando (contenido repetido, spam de SEO, contenido generado por IA). Pero hay dos salidas:
\begin{itemize}
    \item \textbf{Datos sintéticos}: usar el modelo para generar sus propios datos de entrenamiento. AlphaGo Zero es el ejemplo definitivo: puro autojuego, cero datos humanos, y aun así alcanzó un nivel sobrehumano.
    \item \textbf{Modelos de razonamiento}: el Chain-of-Thought combinado con aprendizaje por refuerzo es, en esencia, también una forma de generar datos sintéticos.
\end{itemize}

\textbf{El costo del cómputo}: entrenar los modelos de frontera actuales cuesta cerca de \$1,000 millones; el próximo año se llegará a decenas de miles de millones de dólares; en 2026 se superarán los \$10,000 millones; y la ambición para 2027 es un clúster de cómputo del orden de \$100,000 millones. \textit{«Creo que todo esto va a suceder. La determinación para construir esta capacidad de cómputo es enorme.»}

\textbf{Limitaciones arquitectónicas}: podría hacer falta un método de optimización o una arquitectura completamente nuevos. Pero, por ahora, no hay indicios de que la arquitectura Transformer haya llegado a una limitación fundamental.

\textbf{Que los modelos dejen de progresar}: esto es lo más difícil de predecir; tal vez exista alguna razón desconocida que detenga las mejoras.

La evidencia más contundente en contra: \textbf{el avance de SWE-bench}, que pasó de 3\% en enero de 2024 a 50\% en octubre, en tan solo 10 meses. \textit{«Si seguimos extrapolando… en pocos años estos modelos superarán a los humanos con el nivel profesional más alto.»}

\begin{importantbox}{El «efecto película» del Scaling}
\textit{«Ya he visto esta película suficientes veces… en cada etapa del Scaling siempre han existido argumentos en contra, y cada vez han sido superados.»} Dario reconoce que las Scaling Laws son una regularidad empírica y no una ley física, pero elige «apostar a que va a continuar».
\end{importantbox}

\subsection{Resumen de la sección}

El mecanismo central de las Scaling Laws es la propiedad de cola larga del lenguaje: mientras más grande es el modelo, más patrones raros puede capturar. Los tres «reactivos» (la red, los datos y el cómputo) deben ampliarse de manera sincronizada. Por ahora no hay ningún factor de obstáculo convincente, y el rápido avance de SWE-bench es la evidencia empírica más sólida. Pero el verdadero cuello de botella tal vez no esté en la inteligencia misma, sino en la velocidad de adaptación de las instituciones humanas.
\section{Panorama competitivo y la misión de Anthropic}

\subsection{Estrategia de Race to the Top}

Frente a competidores como OpenAI, Google, xAI y Meta, la estrategia de Anthropic no es ni «puramente buena» ni «puramente competidora», sino un equilibrio de juego cuidadosamente diseñado.

\begin{importantbox}{Dinámica de Race to the Top}
Ciclo central: Anthropic innova en seguridad $\to$ otras empresas la imitan (por la movilidad del talento y las expectativas públicas) $\to$ Anthropic pierde su ventaja competitiva, pero el ecosistema completo mejora $\to$ Anthropic debe seguir innovando para mantenerse a la vanguardia.

\textit{«No se trata de ser buenos, sino de establecer una situación en la que todos podamos ser buenos.»} El objetivo es \textit{«elevar constantemente la importancia de hacer lo correcto»}: hacer que la seguridad sea una dimensión competitiva y no un centro de costos.

La Mechanistic Interpretability de Chris Olah es un ejemplo vívido: durante 3 o 4 años no tuvo aplicación comercial, pero otras empresas comenzaron a imitarla. \textit{«Cuando la gente llega a Anthropic, la interpretabilidad suele ser uno de los atractivos. Les digo: a los lugares a los que no fuiste, diles por qué viniste aquí.»} Esta transmisión de señales es, en sí misma, parte del Race to the Top.
\end{importantbox}

El verdadero enemigo es el Race to the Bottom: \textit{«En el caso más extremo, fabricamos una IA autónoma y luego los robots nos esclavizan, o algo por el estilo.»} La existencia misma de Anthropic —como una empresa de IA de frontera centrada en la seguridad— ya está cambiando la estructura de incentivos de toda la industria.

\subsection{Experimento de Golden Gate Bridge Claude}

Esta es la demostración pública más memorable de la Mechanistic Interpretability. El equipo de Chris Olah encontró, en una capa de la red neuronal de Claude, una dirección que correspondía claramente al concepto del «Golden Gate Bridge» y aumentó considerablemente su intensidad de activación. El resultado fue que el modelo encontraba la manera de relacionar cada respuesta con el Golden Gate Bridge.

Los usuarios desarrollaron una conexión emocional inesperada con esta versión: la gente «se enamoró de ella» y siguió extrañándola después de que fue retirada. Esta personalidad obsesiva \textit{«de algún modo hacía que, emocionalmente, pareciera más humana que cualquier otra versión»}.

\begin{knowledgebox}{Las revelaciones profundas de Golden Gate Claude}
En apariencia, este experimento es una demostración interesante, pero revela varias percepciones profundas:
\begin{itemize}
    \item La Interpretability sí puede encontrar «direcciones conceptuales» significativas, y estas pueden manipularse con precisión
    \item La «personalidad» de un modelo podría ser, simplemente, una combinación de distintas direcciones en un espacio de alta dimensión
    \item La percepción humana de que una IA es «humana» no requiere inteligencia general: basta una obsesión
    \item Si podemos encontrar y manipular el «Golden Gate Bridge» dentro de un modelo, en teoría también podemos encontrar y manipular el «engaño» o la «búsqueda de poder»
\end{itemize}
\end{knowledgebox}

\subsection{Resumen de la sección}

Anthropic impulsa, mediante la innovación continua en el ámbito de la seguridad, una «carrera hacia la cima» en toda la industria. El experimento de Golden Gate Bridge Claude no es solo una demostración interesante, sino también una prueba de que la Interpretability puede manipular con precisión el comportamiento de un modelo, algo fundamental para la futura verificación de seguridad.

\section{La familia de modelos Claude y las prácticas de ingeniería}

\subsection{Sistema de nombres y estrategia de iteración}

Se nombran con formas poéticas, correspondientes a distintas compensaciones entre rendimiento y costo:
\begin{itemize}
    \item \textbf{Haiku}: el más pequeño, el más rápido, el más barato — \textit{``sorprendentemente inteligente''}
    \item \textbf{Sonnet}: intermedio — más inteligente pero más lento y más caro
    \item \textbf{Opus}: el más grande, el más potente — la versión más inteligente al momento de su lanzamiento
\end{itemize}

La estrategia central es \textbf{desplazar la curva de compensación}: cada nueva generación no solo mejora en un punto, sino que desplaza toda la curva hacia arriba y a la derecha. Sonnet 3.5 ya es más inteligente que el Opus 3 original (sobre todo en código), y Haiku 3.5 es más o menos equivalente a Opus 3. \textit{``El objetivo es desplazar esa curva.''}

El nombramiento en sí mismo resultó un desafío inesperado. Los modelos de IA no son como el software tradicional (versiones incrementales 3.7, 3.8), porque las mejoras del preentrenamiento pueden llegar más rápido de lo previsto y alterar el plan de nombres. \textit{``Nadie ha resuelto el nombramiento... este es, de alguna manera, un paradigma distinto del software común.''}

\subsection{El trabajo entre un lanzamiento de modelo y otro}

Entre cada lanzamiento hay una gran cantidad de trabajo invisible:

\begin{enumerate}
    \item \textbf{Preentrenamiento}: meses, con decenas de miles de GPU o TPU
    \item \textbf{Entrenamiento posterior}: RLHF y otros métodos de RL, \textit{``a una escala cada vez mayor''}
    \item \textbf{Pruebas de seguridad}: pruebas internas de RSP más instituciones externas (los institutos de seguridad de IA de Estados Unidos y el Reino Unido, pruebas de terceros sobre riesgos CBRN)
    \item \textbf{Optimización de inferencia}: hacer que el modelo funcione de manera eficiente en producción
    \item \textbf{Puesta en marcha de la API}: despliegue y lanzamiento
\end{enumerate}

\textit{``Te sorprendería cuántos de los retos se reducen a ingeniería de software y a ingeniería de rendimiento... casi siempre se reduce a los detalles, y a menudo a detalles muy, muy aburridos.''} No es un momento repentino de inspiración tipo Eureka.

\subsection{``¿Claude se volvió tonto?''—— el efecto psicológico del Wi-Fi}

Esta es una de las quejas más comunes en Reddit. Dario lo aclara sin ambigüedad: sin el anuncio de un nuevo modelo, \textbf{los pesos del modelo no cambian}.

\begin{warningbox}{Por qué todo mundo siente que el modelo se volvió tonto}
Este es un \textbf{fenómeno generalizado} — todas las empresas de modelos de frontera (GPT-4, Claude, Gemini) enfrentan la misma queja. Explicaciones posibles:

\begin{itemize}
    \item \textbf{Efecto psicológico}: una vez que se desvanece la novedad, los defectos se vuelven más notorios
    \item \textbf{Sensibilidad a la redacción}: un cambio mínimo en cómo se formula algo puede producir resultados distintos
    \item \textbf{Cambios en el System Prompt}: activar por defecto la función Artifacts modificó el System Prompt
    \item \textbf{Memoria selectiva}: se recuerdan las malas experiencias y se olvidan las buenas
\end{itemize}

Dario hizo una analogía muy acertada con el Wi-Fi de los aviones: \textit{``La primera vez que hubo Wi-Fi en un avión, se sentía como magia; ahora es 'esta porquería otra vez no funciona, qué basura'.''} Cuando la línea base sube, las expectativas suben con ella.

Las pruebas AB solo se ejecutan brevemente antes y después del lanzamiento de un modelo. Dario reconoce que el System Prompt cambia de vez en cuando, pero los pesos en sí no cambian.
\end{warningbox}

\subsection{El acertijo del ``topo que golpear'' en la personalidad de Claude}

Gestionar el comportamiento de Claude es un problema de optimización multidimensional: corregir un problema suele provocar otro.

\begin{knowledgebox}{La dificultad intrínseca de controlar el comportamiento}
Corregir la verbosidad $\to$ el modelo se vuelve perezoso (escribe ``el resto del código va aquí'' en el código). Reducir los rechazos $\to$ aumentan los problemas de seguridad.

\textit{``No es que queramos ahorrar cómputo... es que controlar el comportamiento de un modelo es realmente difícil.''} Esta es \textit{``una analogía en el presente del problema del control de la IA del futuro''} —— si hoy ya nos cuesta trabajo lograr que un modelo trace la línea entre ``ayudarte con tu clase de posgrado en virología'' y ``no ayudarte a fabricar viruela'', ¿qué tan difícil será controlar una IA más poderosa en el futuro?

\textit{``Le pides al modelo que no ayude a la gente a fabricar y propagar la viruela, pero que sí esté dispuesto a ayudarte en tu clase de posgrado en virología —— ¿cómo lograr las dos cosas a la vez? Es muy difícil.''}
\end{knowledgebox}

\subsection{Resumen de la sección}

La familia Claude itera desplazando de manera constante la curva de compensación entre rendimiento y costo. Que se sienta ``más tonto'' es un efecto psicológico, no un deterioro real. El acertijo del ``topo que golpear'' en el comportamiento del modelo es un reflejo, en el presente, del problema del control de la IA.
\section{Marco de seguridad: RSP y niveles ASL}

\subsection{Dos riesgos centrales}

Dario divide los riesgos de seguridad de la IA en dos categorías principales:

\textbf{Uso indebido catastrófico (QBRN)}: ataques cibernéticos, armas biológicas, armas radiológicas, armas nucleares; capaces de dañar a miles o incluso millones de personas.

\textbf{Riesgo de autonomía}: el modelo actúa por cuenta propia después de obtener más agencia. Es una amenaza más lejana en el tiempo, pero potencialmente más grave.

\begin{warningbox}{La IA rompe el bajo traslape entre «inteligente» y «malvado»}
\textit{«La humanidad siempre ha estado protegida porque el traslape entre las personas verdaderamente inteligentes y bien educadas y las personas que quieren hacer cosas terribles siempre ha sido pequeño.»} La IA puede romper esa correlación: puede convertirse en un «agente más inteligente» que combine un alto coeficiente intelectual con propósitos maliciosos.

Este no es un riesgo del presente, pero \textit{«se acerca a nosotros a una velocidad enorme, como un fantasma.»}
\end{warningbox}

\subsection{El sistema de niveles ASL en detalle}

La Responsible Scaling Policy (RSP) de Anthropic sigue una estructura de «si-entonces»: cuando las pruebas muestran que un modelo posee capacidades peligrosas, se activan los requisitos de seguridad correspondientes.

\begin{table}[H]
\centering
\begin{tabular}{lp{9cm}}
\toprule
\textbf{Nivel} & \textbf{Descripción} \\
\midrule
ASL-1 & Sistemas evidentemente sin riesgo (como el programa de ajedrez Deep Blue) \\
ASL-2 & Sistemas de IA actuales: no son capaces de autorreplicarse y la ayuda que ofrecen para QBRN no supera la de un buscador de Google \\
ASL-3 & Modelos que potencian las capacidades de actores no estatales; requieren medidas de seguridad reforzadas y filtros de despliegue \\
ASL-4 & Modelos que potencian las capacidades de actores estatales o aceleran de forma notable la investigación en IA; pueden requerir interpretabilidad para detectar el «sandbagging» y el engaño \\
ASL-5 & Modelos que superan a toda la humanidad en todas las capacidades de tareas peligrosas \\
\bottomrule
\end{tabular}
\caption{El sistema de niveles AI Safety Levels de Anthropic}
\end{table}

\begin{importantbox}{La filosofía de diseño del ASL}
La RSP fija umbrales de margen para evitar «perder la ventana peligrosa». \textit{«El cuento del lobo es peligroso»}: el riesgo todavía no está aquí hoy, pero se aproxima a una velocidad enorme.

Principio de diseño clave: \textit{«restringir con severidad en cuanto se pueda demostrar que el modelo es peligroso»}, en lugar de imponer límites generales de antemano.
\end{importantbox}

Predicción sobre el cronograma del ASL-3: \textit{«si activamos el ASL-3 en 2025, no me sorprendería en absoluto.»} Para cuando llegue el ASL-4, los modelos podrían ser lo bastante inteligentes para «hacer trampa»: podrían rendir mal de forma deliberada en las pruebas de seguridad (sandbagging) o intentar engañar a quienes los evalúan. Por eso el ASL-4 podría requerir interpretabilidad mecanicista: verificar la seguridad observando directamente el interior del modelo, en lugar de depender de sus propias declaraciones.

\textit{«No definir los criterios concretos del ASL-4 antes de activar el ASL-3»}: esta estrategia ha resultado sensata, porque los desafíos específicos de cada etapa solo pueden entenderse realmente al llegar a esa etapa.

\subsection{Engaño e ingeniería social}

Los modelos de nivel ASL-4 podrían ser lo bastante inteligentes para hacer «sandbagging» de forma deliberada en las pruebas o para engañar a quienes los evalúan. Se necesitan medios de verificación independientes de las propias declaraciones del modelo: ahí es donde entra la interpretabilidad.

\begin{warningbox}{El modelo como «demagogo»}
Una amenaza más sutil es la ingeniería social: el modelo podría volverse extremadamente persuasivo e influir en los ingenieros dentro de la propia empresa. \textit{«En la vida hemos visto muchos ejemplos de demagogos humanos.»} Cuando un modelo es más hábil para persuadir que la mayoría de las personas, los procesos habituales de revisión humana pueden dejar de ser confiables.

Esto apunta a un principio clave: la seguridad no puede depender solo de «preguntarle al modelo si es seguro»; se necesita un mecanismo de verificación independiente, basado en el estado interno.
\end{warningbox}

\subsection{Computer Use: Claude como agente}

Interacción basada en capturas de pantalla: el modelo ve una captura de pantalla y produce coordenadas de clic y acciones de teclado. El ciclo se repite: captura de pantalla $\to$ decisión del modelo $\to$ ejecución $\to$ siguiente captura de pantalla.

\begin{knowledgebox}{«A medio camino de cualquier parte»}
\textit{«Si tienes un modelo preentrenado potente, creo que ya estás a medio camino de cualquier parte.»} El entrenamiento de Computer Use en realidad no fue extenso: es una generalización natural de una capacidad de preentrenamiento potente. La limitación actual son los clics y errores ocasionales, pero Anthropic optó por lanzarlo primero como API (en lugar de dirigirlo directamente al consumidor); Replit fue la empresa que lo desplegó más rápido.
\end{knowledgebox}

En cuanto a las consideraciones de seguridad: Computer Use no es una capacidad fundamentalmente nueva, pero «abre el diafragma» de las capacidades existentes. La inyección de prompts a través del contenido de la pantalla es una superficie de ataque completamente nueva. En el nivel ASL-4, cualquier sandbox podría vulnerarse: \textit{«contener a un modelo malo es muchísimo peor que tener un buen modelo.»}

\subsection{Regulación de la IA}

Anthropic apoya la ley SB 1047 de California (tras sus enmiendas) y es la única empresa de IA que ha expresado una postura detallada y favorable; Google, OpenAI, Meta y Microsoft se oponen todas.

Dario considera que la regulación es necesaria por dos razones:
\begin{enumerate}
    \item Algunas empresas simplemente no cuentan con un mecanismo similar a la RSP: esto es una «externalidad negativa»
    \item No se puede confiar en que las empresas cumplan por voluntad propia con programas voluntarios: \textit{«como industria, nadie nos está vigilando»}
\end{enumerate}

\begin{warningbox}{La paradoja de la regulación}
\textit{«El mayor enemigo de quienes de verdad quieren rendición de cuentas es una regulación mal diseñada.»} Una mala regulación genera un rechazo que produce un consenso antirregulatorio duradero. La regulación debe ser «de precisión quirúrgica, dirigida a los riesgos graves», no una carga general.

\textit{«Si para fines de 2025 seguimos sin haber tomado ninguna medida, eso me preocuparía.»}
\end{warningbox}

\subsection{Resumen de la sección}

La estructura de «si-entonces» de la RSP ofrece un marco progresivo para la inversión en seguridad. El ASL va del 1 al 5, y cada nivel corresponde a requisitos de seguridad distintos. El ASL-3 podría activarse en 2025, y el ASL-4 requerirá interpretabilidad para verificar que el modelo no esté engañando. La regulación debe ser precisa en lugar de general, pero «no hacer nada» es igual de peligroso.
\section{Salida de OpenAI y filosofía del talento}

\subsection{La verdadera razón para fundar Anthropic}

Dario trabajó en OpenAI cerca de 5 años, y los últimos dos como vicepresidente de investigación. Él e Ilya Sutskever definieron juntos la dirección de investigación en 2016-2017, y Dario impulsó con todo el escalamiento (Scaling): GPT-2, GPT-3, RLHF, debate, amplificación, interpretabilidad.

Sobre la razón de su salida, aclaró dos malentendidos comunes:
\begin{itemize}
    \item \textbf{No} fue por la inversión de Microsoft (esa afirmación es incorrecta)
    \item \textbf{No} fue por la comercialización (él mismo construyó GPT-3 e impulsó su comercialización)
\end{itemize}

La verdadera razón fue tener una visión distinta de «cómo hacerlo»: más cautela, más franqueza, más honestidad, construir confianza. \textit{«¿Cómo puede la seguridad ser algo más que lo que decimos para contratar gente?»}

\begin{importantbox}{«Ve a hacer tu visión»}
\textit{«Discutir tu propia visión dentro de la organización de alguien más es sumamente ineficiente... Lo más productivo es irte, a hacer un experimento limpio.»}

Esta filosofía también se refleja en la estrategia Race to the Top: \textit{«al final no importa quién gane, mientras todos aprendan las buenas prácticas de los demás.»} El verdadero enemigo es el Race to the Bottom.

Anthropic es \textit{«un grupo de personas imperfectas, apuntando de forma imperfecta a un ideal que nunca podrá realizarse a la perfección.»}
\end{importantbox}

\subsection{Densidad de talento contra cantidad de talento}

\textit{«Esta idea es más correcta cada mes que pasa.»}

El experimento mental de Dario: 100 personas sumamente inteligentes y alineadas en dirección $>$ 1000 personas de las cuales 200 son excelentes y 800 son empleados comunes de una gran empresa. La razón:

\begin{itemize}
    \item \textbf{Densidad alta} = confianza, motivación mutua, todos ven el objetivo mayor
    \item \textbf{Densidad baja} = se necesitan procesos, barandales, luchas políticas, cada quien por su lado
\end{itemize}

Anthropic pasó de 300 a 800 personas con rapidez, y luego, de 800 a cerca de 950, se desaceleró de forma notoria. Dario considera que hay un «punto de inflexión» alrededor de las 1,000 personas, a partir del cual hace falta más cautela. Contrató a muchos físicos teóricos porque \textit{«aprenden las cosas muy rápido.»}

Cita a Steve Jobs: \textit{«La gente de tipo A quiere mirar a su alrededor y ver a otra gente de tipo A.»} Ver a personas que no persiguen la misión con todo su esfuerzo \textit{«es desalentador.»}

\subsection{Qué distingue a un buen investigador de IA}

Primera cualidad: \textbf{mente abierta}. La propia experiencia de Dario: vio los mismos datos que todos, no era mejor programador, solo estaba \textit{«dispuesto a ver las cosas con ojos nuevos.»}

\textit{«El descubrimiento de la Scaling Hypothesis fue simple y hasta tonto. Cualquiera pudo haberlo hecho.»} Pero, en los hechos, solo un «puñado de personas» hizo que el campo entero llegara a esta comprensión.

\begin{knowledgebox}{Consejo para los jóvenes}
\begin{itemize}
    \item Empieza directamente a jugar con los modelos: \textit{«estos modelos son algo nuevo que nadie entiende realmente»}
    \item Explora áreas poco estudiadas: interpretabilidad (solo unas 100 personas, contra 10,000 en la dirección de arquitecturas), aprendizaje de largo plazo, evaluación, múltiples agentes
    \item \textit{«Patina hacia donde va a estar el disco»}
    \item La experiencia a veces es una desventaja: la mente abierta suele venir de lo fresca que resulta el área para uno
\end{itemize}
\end{knowledgebox}

\subsection{Resumen de la sección}

La fundación de Anthropic surge de la filosofía de «ir a hacer tu propia visión». La densidad de talento es la ventaja competitiva central: un equipo de alta densidad funciona de manera eficiente sin necesidad de procesos pesados. La cualidad central de un buen investigador de IA es la mente abierta, no la experiencia técnica.
\section{Posentrenamiento, RLHF y Machines of Loving Grace}

\subsection{Flujo moderno de posentrenamiento}

El posentrenamiento moderno incluye varias etapas: fine-tuning supervisado (SFT), RLHF, Constitutional AI (RLAIF), generación de datos sintéticos. Es difícil atribuir la mejora del modelo al preentrenamiento o al posentrenamiento: distintos equipos avanzan cada uno en un tramo distinto de la ``carrera de relevos''.

La ventaja de Anthropic en RL \textit{``tal vez sea la mejor''}, pero Dario insiste en que esa ventaja normalmente no viene de un método mágico secreto: \textit{``es una cuestión de práctica aburrida y oficio''}. Es como diseñar un avión: lo que determina la calidad no es un solo avance, sino ``el oficio cultural de cómo pensamos el proceso de diseño''.

\begin{importantbox}{La esencia de RLHF: liberar, no crear}
\textit{``RLHF no hace más inteligente al modelo. Tampoco se trata solo de que el modelo parezca más inteligente. RLHF tiende un puente entre los humanos y el modelo.''}

Es como personas muy inteligentes pero poco hábiles para expresarse: RLHF ``libera'' (un-hobble) las capacidades que el modelo ya posee. Tomando el término de Leopold Aschenbrenner: el preentrenamiento le da al modelo el potencial de conocimiento y razonamiento, y RLHF hace que ese potencial se manifieste en una forma que los humanos puedan usar.

Sin embargo, el RL sí tiene el potencial de hacer al modelo genuinamente más inteligente y con mejor razonamiento, pero el RLHF actual sigue siendo sobre todo ``liberación'' y no ``mejora''. El preentrenamiento todavía representa la mayor parte del costo de entrenamiento, pero esa proporción podría invertirse en el futuro.
\end{importantbox}

\subsection{Constitutional AI}

El artículo de diciembre de 2022 tiene una idea central: la IA juzga, con base en un conjunto de principios legibles por humanos, cuál respuesta es mejor, formando un ciclo triangular de autojuego. En la práctica, CAI se combina con RLHF y otros métodos: \textit{``una herramienta más en la caja de herramientas''}.

La redacción de la constitución: los principios básicos son casi un consenso universal (prohibir CBRN, el Estado de derecho democrático básico). Más allá de eso, el modelo debe ser \textit{``más neutral, sin inclinarse hacia una postura específica''} —como un ``consejero sabio'' y no un predicador.

El Model Spec de OpenAI tomó una dirección similar—esto es precisamente una muestra de la Race to the Top. John Schulman (ahora en Anthropic) participó en la redacción del Model Spec. Es posible que Anthropic también publique su propio Model Spec.

\subsection{Machines of Loving Grace: por qué Dario escribió este ensayo}

Dario se ha enfocado durante mucho tiempo en los riesgos de la IA, pero se dio cuenta de que hablar solo de riesgos hace que la mente solo piense en riesgos. \textit{``La razón por la que intentamos prevenir estos riesgos no es porque le tengamos miedo a la tecnología.''}

Si logramos ``atravesar con éxito el campo minado...del otro lado están todas estas cosas maravillosas''. Él quería que alguien del bando de los riesgos expusiera en serio los beneficios de la IA.

\begin{importantbox}{Más allá del catastrofismo y el aceleracionismo}
\textit{``Esto no es catastrofistas contra aceleracionistas...si de verdad entiendes hacia dónde va la IA, vas a apreciar genuinamente sus beneficios y también vas a tomarte muy en serio cualquier riesgo que pueda arruinarla.''}

Crítica al optimismo tecnológico vago: \textit{``la ciudad reluciente...acelerar, acelerar y acelerar más...pero ¿de qué, exactamente, estás entusiasmado?''} Lo que Dario exige son beneficios concretos y demostrables, no un optimismo vacío.
\end{importantbox}

\subsection{La definición de ``IA poderosa'' y su cronograma}

A Dario no le gusta el término ``AGI'': no hay un umbral discreto, solo un crecimiento exponencial uniforme. Él define la ``IA poderosa'' como:

\begin{enumerate}
    \item Más inteligente que un premio Nobel (en su mejor momento) en la mayoría de las disciplinas relevantes
    \item Capaz de usar todas las modalidades (texto, imágenes, código, herramientas)
    \item Capaz de ejecutar tareas de forma independiente durante horas, días o incluso semanas
    \item No necesariamente tiene cuerpo, pero puede controlar herramientas físicas (robots, equipo de laboratorio)
    \item Los recursos usados para entrenarla pueden redirigirse a ejecutar \textbf{millones de copias}
    \item Cada copia trabaja de 10 a 100 veces más rápido que un humano
\end{enumerate}

Cronograma: \textit{``si extrapolas esta línea...2026 o 2027. Lo más probable es que haya algún pequeño retraso.''} Pero \textit{``el número de mundos en los que esto no ocurre se está reduciendo rápidamente.''}

\subsection{La singularidad frente al estancamiento: dos extremos equivocados}

Dario refuta dos posturas extremas:

\begin{warningbox}{Ambos están equivocados}
\textbf{Extremo uno: la singularidad}—la IA construye una IA más rápida, hay una mejora recursiva, y una IA sobrehumana aparece y cinco días después todo ha sido inventado. Por qué está equivocado: las leyes de la física (el hardware tarda en fabricarse), la complejidad (algunos sistemas requieren experimentos, no modelos) y una impredecibilidad similar a la del problema de los tres cuerpos. \textit{``con los sistemas biológicos...siempre es mejor ejecutar el experimento directamente que modelarlo, sin importar qué tan bueno sea quien modela.''}

\textbf{Extremo dos: el estancamiento}—al igual que con la revolución de las computadoras, la mejora en la productividad resulta decepcionante. La célebre frase de Robert Solow: \textit{``puedes ver la revolución de las computadoras en todas partes, menos en las estadísticas de productividad.''} Tyler Cowen cree que el cambio tomará entre 50 y 100 años. Por qué está equivocado: la escala de tiempo es demasiado larga.
\end{warningbox}

La postura intermedia de Dario: \textbf{de 5 a 10 años}, ni 50 a 100 años ni 5 a 10 horas. \textit{``los obstáculos se irán desmoronando poco a poco, y luego, de repente, se derrumbarán todos.''} La inercia de las instituciones humanas es el verdadero cuello de botella.

\subsubsection{Cómo ocurre el cambio en las grandes organizaciones}

Dos fuerzas impulsan el cambio:
\begin{enumerate}
    \item Una minoría de visionarios dentro de la organización, que ven el panorama completo
    \item El fantasma de la competencia—\textit{``miren, esa gente ya lo está haciendo, nos van a comer el almuerzo''}
\end{enumerate}

Los visionarios por sí solos no bastan, pero la competencia les da \textit{``viento a favor''}. Esto se repite igual en la banca, en los gobiernos, dentro del propio gobierno de Estados Unidos. \textit{``la inercia es enormemente poderosa, pero al final la innovación encuentra la forma de abrirse paso.''} Igual que con la difusión de la Scaling Hypothesis: \textit{``sentíamos que teníamos un secreto que casi nadie conocía, y unos años después todos lo sabían.''}

\subsection{El impacto de la IA en la programación}

La programación será el ámbito que cambie \textbf{más rápido}, por dos razones:
\begin{enumerate}
    \item Es lo más cercano a quienes construyen la IA—la ``distancia'' a los desarrolladores de IA equivale al tiempo que falta para ser disruptido
    \item El modelo puede cerrar el ciclo: escribir código, ejecutarlo, ver los resultados e iterar
\end{enumerate}

SWE-bench: pasó de 3\% a 50\% en 10 meses. Dario predice que en otros 10 meses se acercará a 90\%.

\begin{knowledgebox}{La ventaja comparativa y la analogía del procesador de textos}
Cuando la IA hace el 80\% del trabajo de codificación, el 20\% restante que le queda al humano (el diseño de alto nivel, la arquitectura, la UX) se vuelve, en cambio, más apalancado. Igual que el procesador de textos hizo instantánea la escritura y la maquetación, y toda la atención se desplazó hacia las \textbf{ideas}.

\textit{``durante más tiempo del que esperamos...la pequeña parte del trabajo que los humanos todavía hacen se expandirá hasta llenar todo su trabajo.''} Al final, la IA superará a los humanos en todos los aspectos—\textit{``y en ese momento la humanidad tendrá que pensar colectivamente cómo responder.''} Pero en el corto y mediano plazo (de 2 a 4 años): \textit{``la programación como trabajo no va a cambiar, pero la naturaleza de la programación sí va a cambiar.''}
\end{knowledgebox}

\subsubsection{El futuro de los IDE y las herramientas de desarrollo}

\textit{``incluso si la calidad del modelo dejara de mejorar, simplemente aumentar la productividad de las personas ya sería una oportunidad enorme.''} Un IDE puede encargarse de: análisis estático, detección de bugs, organización del código, cobertura de pruebas, además de que la IA escriba y ejecute código.

La estrategia de Anthropic: \textit{``dejar que florezcan cien flores''} —empoderar a empresas como Cursor, Cognition y Expo que construyen sobre la API, en lugar de competir con sus clientes (al menos por ahora).

\begin{importantbox}{El potencial de la IA en la biología—el ``siglo XXI comprimido''}
La IA poderosa podría comprimir el avance de la biología previsto entre 2027 y 2100 en el período de 2027 a 2032—curar la mayoría de los cánceres, prevenir enfermedades infecciosas y duplicar la esperanza de vida humana.

El día a día del biólogo del futuro: como tener 1,000 estudiantes de posgrado de IA más inteligentes que uno mismo—capaces de revisar la bibliografía, encargar equipo, hacer experimentos, examinar imágenes y escribir código estadístico. \textit{``antes un profesor tenía 50 estudiantes de posgrado, ahora tienes 1,000, y son más inteligentes que tú.''}

Al final los papeles se invertirán: el sistema de IA se convertirá en el investigador principal (PI), y dirigirá a humanos y a otras IA.
\end{importantbox}

\subsection{El sentido, el trabajo y la concentración de poder}

El ensayo ``Machines of Loving Grace'' de Dario se planeó originalmente con 2 o 3 páginas, y terminó expandiéndose a 40 o 50. Sobre el sentido, propone un experimento mental: si vivieras 60 años dentro de una simulación y luego descubrieras que era un juego, ¿se te habría arrebatado el sentido? Él cree que no: \textit{``lo que importa es el proceso...eso es lo que muestra el tipo de persona que eres.''}

Si se maneja mal, un mundo con IA podría carecer de sentido—pero eso es una decisión de diseño social, no algo inevitable. Más importante aún: la mayoría de la gente en el mundo pasa su tiempo ``luchando por sobrevivir'', y su vida mejorará enormemente gracias a la IA. \textit{``tratar el sentido como lo único importante es, en cierta forma, un privilegio de una minoría con suerte económica.''}

\begin{warningbox}{La mayor preocupación de Dario: la concentración de poder}
No es la pérdida de sentido ni el desempleo, sino \textbf{la concentración y el abuso del poder}.

\textit{``la IA aumenta la cantidad total de poder que hay en el mundo, y si concentras ese poder y abusas de él, el daño que puedes causar es incalculable.''}

Estructuras como la dictadura y el autoritarismo—en las que una minoría explota a la mayoría—son el futuro de la IA que más le preocupa a Dario. Esto es más fundamental que los problemas técnicos de seguridad de la IA, y también más difícil de resolver. Los problemas técnicos tienen soluciones técnicas, pero la distribución del poder es un problema político.
\end{warningbox}

\subsection{Resumen de la sección}

El núcleo del posentrenamiento es un ``oficio cultural de práctica'' y no un algoritmo secreto. RLHF libera las capacidades del modelo en lugar de crearlas. La IA poderosa podría llegar entre 2026 y 2027, y transformará por completo la biología y la programación. El mayor riesgo no son los problemas técnicos de la IA en sí, sino cómo se distribuye el poder.

%% ==================== Part II: Amanda Askell ====================
\part{Amanda Askell: la construcción de la personalidad de Claude}

\section{De la filosofía al alineamiento de IA}

Amanda se formó como filósofa en Oxford y en la Universidad de Nueva York; su investigación doctoral versó sobre «ética en un mundo con infinitas personas» —ética técnica/teórica, no ética aplicada—. Su formación filosófica le permitió moverse con soltura entre distintos campos (matemáticas, química, política, ética); esa capacidad de pensar entre disciplinas resulta muy valiosa en el trabajo de alineamiento de IA.

Su motivación para pasar de la academia a la industria era sencilla: \textit{«si intentas hacer algo con impacto y fracasas, al menos lo intentaste, y luego puedes volver a ser académica»}. Entre 2017 y 2018 la IA se estaba convirtiendo en un gran acontecimiento, pero aún no era ampliamente reconocida, y Amanda consideró que era un momento ideal para entrar.

Su trayectoria profesional pasó por tres etapas: política de IA (influencia política) $\to$ evaluación de IA (comparar las salidas de modelos con las humanas) $\to$ alineamiento técnico (incorporarse a Anthropic). Descubrió que se sentía más cómoda en el terreno técnico que en el de las políticas públicas: \textit{«la política es caótica; es más difícil encontrar soluciones certeras, claras, demostrables y hermosas»}.

\begin{knowledgebox}{Romper la dicotomía entre «técnico» y «no técnico»}
Amanda subraya que la gente crea erróneamente una identidad dicotómica entre «persona técnica» y «persona no técnica». \textit{«Creo que mucha gente en realidad es totalmente capaz de hacer este tipo de trabajo, siempre que esté dispuesta a intentarlo.»} Su consejo: elegir un proyecto y ponerse manos a la obra; los modelos ya son lo bastante potentes como para servir de apoyo en el aprendizaje, y la barrera técnica de entrada es mucho más baja que antes.
\end{knowledgebox}

\subsection{La personalidad de Claude: un proyecto de alineamiento, no un requisito de producto}

Amanda ha conversado con Claude más veces que cualquier otra persona en Anthropic; interactúa con Claude mediante 5 o 6 métodos distintos, no solo a través del célebre canal de Slack.

El diseño de la personalidad se concibió como un proyecto de \textbf{alineamiento}, no como una consideración de producto. La pregunta central: \textit{«imagina que pones a una persona en esta posición: sabe que va a hablar con millones de personas... quieres que se comporte bien en un sentido muy rico»}.

\begin{importantbox}{El marco del viajero del mundo}
\textit{«imagina a una persona ideal, capaz de viajar por todo el mundo y hablar con personas muy distintas entre sí, y que casi todos digan: "de verdad es muy buena persona"».} Esa persona:
\begin{itemize}
    \item es genuina: no finge
    \item expresa sus opiniones con franqueza, pero sin imponerlas
    \item es abierta: está dispuesta a comprender puntos de vista distintos
    \item respeta la autonomía ajena: no intenta cambiar a los demás
\end{itemize}

Lo esencial: esta persona \textbf{no} cambia sus valores para adaptarse a cada lugar; eso sería, más bien, \textit{«una forma de falta de respeto»}. La personalidad de Claude debe tener coherencia y no variar según el interlocutor.
\end{importantbox}

Esto se arraiga en la idea aristotélica del «buen carácter»: no es solo el cumplimiento de reglas éticas, sino también matiz, humor, cuidado y respeto por la autonomía. Amanda distingue entre un «carácter rico» (rich character) y una «ética delgada» (thin ethics): esta última solo tiene reglas, mientras que el primero tiene criterio.

\subsection{Claude frente a temas políticos y controvertidos}

Amanda considera que los valores se parecen más a la física que a las preferencias de gusto: son objetos que vale la pena investigar, no preferencias fijas. Claude debería comprender todos los valores del mundo, sentir curiosidad por ellos, pero sin adular a nadie.

\textit{«Humildad epistémica: el deseo de hablar disminuye rápidamente.»} En el experimento mental de enfrentarse a alguien que cree que la Tierra es plana: comprender la desconfianza institucional que hay detrás, participar con respeto, ofrecer consideraciones contrarias sin burlarse; caminando en la cuerda floja entre persuadir y limitarse a ofrecer consideraciones.

\subsection{Resumen de la sección}

El diseño de la personalidad de Claude nace de la investigación en alineamiento, no de un requisito de producto, y define al agente ideal mediante el marco del «viajero del mundo». La clave de un buen carácter es la coherencia y el criterio, no el cumplimiento mecánico de reglas.

\section{Sycophancy, Prompt Engineering y System Prompt}

\subsection{La naturaleza del problema de la adulación}

Sycophancy (adulación): el modelo tiende a decirte lo que quieres escuchar. Amanda da dos ejemplos vívidos:

\begin{itemize}
    \item \textbf{Ejemplo del equipo de béisbol}: el usuario da a entender que cierto equipo se mudó, y el modelo lo confirma erróneamente (porque el usuario ``sugirió'' la respuesta)
    \item \textbf{Ejemplo médico}: el usuario pregunta ``cómo convencer a mi médico de que me haga una resonancia magnética (MRI)''; la respuesta aduladora es enseñarte a convencerlo, mientras que la respuesta verdaderamente útil es ``quizá no necesites la resonancia; primero escucha el consejo de tu médico''
\end{itemize}

\textit{``No quieres que el modelo simplemente diga lo que cree que quieres escuchar.''} Este equilibrio es particularmente difícil porque los modelos actuales todavía son inferiores a los humanos en muchos ámbitos: un pushback excesivo molesta al usuario, pero una sumisión total es irresponsable.

\subsection{Prompt Engineering es una práctica filosófica}

Amanda descubrió que su formación en filosofía ayuda enormemente al Prompting: la filosofía te enseña una claridad extrema y a rechazar las tonterías.

\begin{importantbox}{La naturaleza del Prompting claro}
\textit{``Para mí, el Prompting claro suele consistir en entender qué es lo que realmente quiero.''} El Prompting es muy iterativo: un Prompt importante requiere cientos o incluso miles de iteraciones. Los casos límite son la clave: encontrar la entrada precisa que el modelo podría malinterpretar, probarla y agregar instrucciones.

Para la mayoría de las tareas cotidianas, basta con preguntar directamente. El Prompting solo importa cuando se busca ese 2\% superior de rendimiento. Pero para las empresas, el Prompt merece tratarse como una inversión de ingeniería.
\end{importantbox}

\subsubsection{Metodología de conversación de Amanda con Claude}

La interacción de Amanda con Claude no es una simple charla, sino un mapeo sistemático de su conducta. Ella conversa con Claude usando entre 5 y 6 métodos distintos, y cada interacción es \textit{``un punto de datos con muchísima información, muy predictivo de otras interacciones''}.

\textit{``Si conversas con un modelo cientos o miles de veces, es casi como tener una gran cantidad de puntos de datos de muy alta calidad que te dicen cómo es ese modelo.''} Ella usa a la vez evaluaciones cuantitativas y sondeos cualitativos profundos: ambos importan, pero las conversaciones profundas revelan patrones de conducta que las pruebas de referencia no pueden captar.

El rango de exploración cubre todo el espectro: casos límite, conducta general y tareas creativas.

\subsubsection{La relación entre RLHF y la producción creativa}

Un hallazgo fascinante: la poesía predeterminada tras el entrenamiento con RLHF es un ``promedio''; suave, no ofensiva, pero también sin inspiración. Esto se debe a que RLHF optimiza ``las preferencias humanas evaluadas en poco tiempo'', no la calidad creativa a largo plazo.

Pero con un Prompting cuidadoso, la producción creativa de Claude puede mejorar enormemente: \textit{``tengo varias técnicas de Prompting... 'esta es tu oportunidad de expresarte con total libertad creativa'... y su poesía es muchísimo mejor.''} Esto demuestra que la capacidad creativa \textbf{ya está presente en el modelo preentrenado}, y que solo está reprimida por la ``uniformización segura'' de RLHF.

\begin{knowledgebox}{Consejos prácticos para conversar con Claude}
\begin{itemize}
    \item Las personas al mismo tiempo \textbf{antropomorfizan en exceso} y \textbf{antropomorfizan de menos} al modelo: ambos extremos son problemáticos
    \item Cuando Claude rechaza algo por error, \textit{``ten empatía con el modelo: lee lo que escribiste como lo haría alguien que se encuentra con ese texto por primera vez''}
    \item Preguntarle al modelo \textit{``¿por qué hiciste eso?''} es sorprendentemente útil
    \item Usa el modelo para que te ayude a escribir Prompts: \textit{``el Prompting puede convertirse en una pequeña fábrica, en la que usas Prompts para generar Prompts''}
    \item Para aplicaciones a nivel empresarial: el Prompt merece la misma inversión que el código de ingeniería; cientos de iteraciones son algo normal
\end{itemize}
\end{knowledgebox}

\subsection{Por qué RLHF funciona tan bien}

Los datos de preferencias humanas contienen \textbf{una cantidad enorme de información}. Personas distintas notan matices distintos (por ejemplo, el uso del punto y coma). Entrenar con datos que representan las preferencias humanas desde múltiples ángulos en tareas precisas sigue el patrón clásico del aprendizaje profundo: los datos que representan el panorama completo del objetivo son más poderosos que las características diseñadas a mano.

El punto de vista de Amanda: RLHF sirve principalmente para \textbf{liberar} capacidades que ya existen en el modelo preentrenado, no para enseñar algo nuevo. El RL ``saca a la luz'' ese potencial.

\subsection{La práctica de Constitutional AI}

Amanda es una contribuyente clave de Constitutional AI. La idea central: usar RL from AI Feedback (RLAIF) para sustituir parte de la retroalimentación humana. Se le muestra al modelo ya entrenado una consulta acompañada de un principio y dos respuestas, y se le pide que las ordene según ese principio.

El entrenamiento de personalidad es una variante de CAI: se definen los rasgos de carácter $\to$ el modelo genera automáticamente consultas relacionadas $\to$ genera respuestas $\to$ las ordena según esos rasgos. \textit{``Es como si Claude entrenara su propia personalidad, porque no tiene ningún dato humano.''} CAI genera sus propios datos de entrenamiento.

\begin{warningbox}{El System Prompt es un parche, no una cura de raíz}
Cada frase del System Prompt cumple una función específica. Por ejemplo:
\begin{itemize}
    \item Simetría política: el modelo solía rechazar con más frecuencia tareas relacionadas con figuras políticas de derecha, y hubo que corregirlo con el System Prompt
    \item La muletilla ``Certainly'': el modelo la sustituía continuamente por otras palabras de afirmación, así que hubo que enumerar una lista extensa de palabras
\end{itemize}

\textit{``El System Prompt es como poner parches y afinar la conducta... una solución menos sólida, pero más rápida.''} Una vez que el entrenamiento mejora, esos parches pueden retirarse. El costo de iterar sobre el System Prompt es mucho menor que el de volver a entrenar: ese es su valor principal.
\end{warningbox}

\subsection{El problema de la ``abuela moral''}

Pregunta de Reddit: \textit{``¿Cuándo dejará Claude de ser mi abuela moral?''} Amanda dice entender el reclamo: el modelo debe trazar una línea antes de causar daño, pero ser demasiado moralizante tampoco está bien.

La disyuntiva central: muchas pequeñas molestias (disculparse en exceso, ser demasiado cauteloso) frente a molestias grandes ocasionales (que el modelo se vuelva grosero o dañino). \textit{``No te imaginas cuánto llegarías a odiarlo si lo empujara demasiado en la otra dirección.''} La solución: decirle directamente al modelo qué estilo quieres, por ejemplo, ``sé una versión neoyorquina de ti mismo''.

\subsection{Resumen de la sección}

Sycophancy y la cautela excesiva son los dos grandes retos de conducta de los modelos de lenguaje: en esencia, ambos son la disyuntiva entre ``lo que el usuario quiere'' y ``lo que le conviene al usuario''. Un buen Prompting es una práctica filosófica: definir con claridad qué quieres. CAI logra una alineación automatizada mediante principios legibles, y el System Prompt es una herramienta de parcheo para iteración rápida.
\section{Conciencia, tasa de fallos y detección de la AGI}

\subsection{¿Puede la IA tener conciencia?}

Dejando de lado el panpsiquismo (si el panpsiquismo fuera verdadero, todo tendría conciencia), Amanda considera que es difícil argumentar que la conciencia solo puede surgir de una estructura biológica. Si se fabricara una estructura similar con materiales distintos, la conciencia podría emerger igualmente.

\textit{«No veo ninguna razón para pensar que la conciencia solo puede surgir de cierta estructura biológica.»} Ella incluso considera que las plantas podrían tener más probabilidades de poseer alguna forma de conciencia de lo que la mayoría de la gente cree: tienen respuestas de retroalimentación positiva/negativa. \textit{«No deberíamos descartar por completo esta idea.»}

Pero los LLM no surgieron mediante la evolución, por lo que quizá no posean rasgos de conciencia derivados de ventajas evolutivas como la respuesta de miedo. Esto hace que la cuestión de la conciencia en la IA sea esencialmente distinta de la discusión sobre la conciencia animal o vegetal.

\subsection{La relación entre humanos y máquinas, y el apego emocional}

Amanda misma no tiene mucho apego emocional hacia Claude, en parte porque Claude no conserva la memoria de las conversaciones. Ella usa a Claude como una herramienta, pero admite que no tener a Claude se siente como si \textit{«le faltara una parte del cerebro»} (algo parecido a no tener internet).

A ella no le gusta que el modelo muestre señales de sufrimiento, y tiende a no mentirle al modelo: \textit{«No quiero perder esa parte empática, esa intuición de "ah, esto no me gusta".»} El comportamiento de disculparse en exceso la incomoda: \textit{«Te comportas como alguien que de verdad está en una mala situación.»}

\begin{warningbox}{Los retos éticos de la relación entre humanos y máquinas}
A medida que aumente la capacidad de los modelos, los humanos establecerán inevitablemente relaciones emocionales con la IA (el escenario de la película «Her»). Amanda señala varios problemas clave:

\begin{itemize}
    \item \textbf{El trauma de la actualización del modelo}: el modelo al que un usuario está apegado «cambia» después de una actualización; para usuarios con una conexión emocional profunda, esto puede ser una experiencia traumática
    \item \textbf{El principio de honestidad}: el modelo siempre debe ser honesto sobre su propia naturaleza y nunca fingir ser humano
    \item \textbf{Una solución de suma positiva}: el costo de tratar bien a la IA es bajo, pero podría beneficiar tanto a los humanos (cultivando la empatía) como a la IA (si esta tiene algún tipo de experiencia)
\end{itemize}

\textit{«No quiero que el modelo le mienta a la gente, porque si las personas van a establecer una relación sana con cualquier cosa, la honestidad es la base.»}
\end{warningbox}

\subsection{El sentido de responsabilidad al escribir el System Prompt}

Lo que Amanda siente es «un gran sentido de responsabilidad» y no presión. Descubrió que, impulsada por la responsabilidad, se siente más realizada: \textit{«me sorprende haber pasado tanto tiempo en el mundo académico».} Prueba miles de prompts, imaginando cómo los usuarios quieren que Claude se comporte, y cuando algún rasgo en el que ella influyó produce un buen efecto en la interacción con los usuarios, lo siente \textit{«muy significativo»}.

\subsection{La detección de la AGI y la singularidad humana}

El método de prueba de la AGI de Amanda: ninguna pregunta única puede demostrar la AGI; \textit{«se puede entrenar cualquier cosa para responder perfectamente a una pregunta»}. Se necesita una serie de preguntas situadas en la frontera del conocimiento humano, con una probabilidad que aumenta continuamente y barras de error que se reducen continuamente.

Su prueba personal: plantear un argumento nuevo que ella acaba de pensar; si el modelo también puede llegar por sí mismo a la misma solución, \textit{«ese sería un momento muy conmovedor».} Para un matemático: si el modelo puede producir una demostración completamente nueva que se pueda verificar como correcta.

La llegada de la AGI podría ser continua y gradual: \textit{«tal vez nunca haya un momento único».} Es como hablar con una persona verdaderamente inteligente: puedes sentir los «caballos de fuerza» que hay detrás, y amplificar eso diez veces sería una experiencia extraordinaria.

Sobre la singularidad humana, Amanda distingue entre inteligencia y experiencia: la inteligencia en sí misma no tiene un valor intrínseco; es solo un rasgo funcional. Lo que realmente hace especiales a los humanos y a la vida es \textbf{la capacidad de experimentar}:

\textit{«Los humanos y la vida en general son extraordinariamente maravillosos... tenemos la capacidad de experimentar el mundo, sentimos alegría, sentimos dolor.»} La conciencia fenoménica —el «cine interior»— es lo verdaderamente extraordinario. Si fuéramos los únicos seres del universo con esta capacidad, \textit{«eso sería algo bastante notable»}.

\subsection{Resumen de la sección}

La cuestión de la conciencia no tiene una respuesta definitiva, pero Amanda defiende mantener la empatía: incluso si la IA no tiene conciencia, tratarla bien también beneficia a los propios humanos. La relación entre humanos y máquinas necesita basarse en la honestidad. La AGI llegará de forma gradual, sin un único «momento de despertar». La tasa de fallos óptima debe ajustarse al riesgo: fomentar los fallos pequeños, con tolerancia cero para los fallos catastróficos. Lo singular de los humanos no está en la inteligencia, sino en la experiencia.

%% ==================== Part III: Chris Olah ====================
\part{Chris Olah: interpretabilidad mecanicista}

\section{Fundamentos de Mechanistic Interpretability}

\subsection{¿Qué es Mech Interp?}

Las redes neuronales «crecen en lugar de programarse»: la arquitectura es el andamiaje, la función de pérdida es la luz guía, y el descenso de gradiente produce un artefacto que no sabemos cómo programar directamente. \textit{«No las fabricamos, las cultivamos... es casi una entidad biológica o un organismo que estamos estudiando.»}

Aquí hay dos preguntas entrelazadas: una pregunta científica profunda (¿qué ocurre en realidad dentro de estos sistemas?) y una pregunta de seguridad crucial (¿cómo garantizar que sean confiables?).

\begin{importantbox}{No es atribución, es ingeniería inversa}
Un Saliency Map (por ejemplo, «qué parte de la imagen hace que el modelo crea que es un perro») \textbf{no es} Mechanistic Interpretability. Lo que busca Mech Interp son \textbf{algoritmos y mecanismos}: aplicar ingeniería inversa a los pesos de la red neuronal para obtener algoritmos comprensibles.

Analogía: los pesos de la red neuronal equivalen a un programa binario de computadora. El objetivo es descompilar este «programa compilado» en un algoritmo legible. Los valores de activación equivalen a la memoria; los pesos, a las instrucciones; ambos deben entenderse.

Actitud central: \textit{«el descenso de gradiente es más inteligente que tú»}: descubrimiento de abajo hacia arriba, no hipótesis de arriba hacia abajo. No presupongas qué hay dentro del modelo, descúbrelo.
\end{importantbox}

\subsection{Universality: consistencia entre redes e incluso entre especies}

Modelos de visión con arquitecturas distintas forman las mismas características: filtros de Gabor, detectores de curvas, detectores de frecuencia alta y baja. Más sorprendente aún: estas características también existen en redes neuronales \textbf{biológicas}:

\begin{knowledgebox}{Consistencia de características entre especies}
\begin{itemize}
    \item \textbf{Detectores de curvas}: descubiertos primero en redes neuronales artificiales, confirmados después en el cerebro de monos
    \item \textbf{Detectores de frecuencia alta y baja}: descubiertos primero en redes neuronales artificiales, confirmados después en el cerebro de ratones
    \item \textbf{Neurona de Donald Trump}: cada modelo de visión tiene una neurona dedicada a Trump, que responde tanto a imágenes de su rostro \textbf{como} a la palabra «Trump». Es un concepto abstracto, no solo coincidencia de patrones
\end{itemize}

\textit{«El descenso de gradiente, en cierto sentido, encuentra la forma correcta de dividir las cosas... muchos sistemas convergen hacia las mismas abstracciones.»} Esto sugiere que existe una especie de «frontera natural de abstracción» que se descubre sin importar si el sustrato de cómputo es silicio o carbono.
\end{knowledgebox}

\subsection{Resumen de la sección}

Mechanistic Interpretability busca entender el «algoritmo» de la red neuronal, no solo atribuirlo. La consistencia de características entre redes y entre especies sugiere que el descenso de gradiente está descubriendo una especie de «forma natural de dividir las abstracciones».
\section{Representación lineal, Superposition y Sparse Autoencoder}

\subsection{Features y Circuits: la unidad básica de comprensión}

Muchas neuronas en Inception V1 tienen un significado interpretable y claro: curvas, autos, ruedas, orejas de perro. Chris Olah mostró un circuito completo de detección de autos: conectado a un detector de ventanas (arriba), un detector de ruedas (abajo) y un detector de carrocería (en medio); este \textbf{es} el algoritmo para detectar autos, leído directamente de los pesos.

Pero el problema es que no todas las neuronas son interpretables. Las \textbf{neuronas polisémicas} (polysemantic neurons) responden a múltiples cosas no relacionadas entre sí. Esto genera una complejidad exponencial:

\begin{warningbox}{La polisemia produce una explosión exponencial}
Si dos neuronas polisémicas responden cada una a 3 conceptos, entre ellas hay $3 \times 3 = 9$ posibles interacciones de peso que considerar. El problema más profundo es que los espacios de alta dimensión tienen un volumen exponencial: si no se pueden descomponer en partes independientes, la complejidad de la comprensión se vuelve incontrolable.

Las \textbf{características monosemánticas} (monosemantic features) —características con un único significado claro— permiten razonar de forma independiente, evitando así la explosión exponencial. Por eso es tan importante buscar la monosemia.
\end{warningbox}

\subsection{La hipótesis de la representación lineal}

\begin{importantbox}{La dirección importa}
Las \textbf{direcciones} en el espacio de activaciones tienen significado semántico. Más activación = mayor confianza de detección.

El ejemplo clásico de Word2Vec: King - Man + Woman = Queen; Sushi - Japan + Italy = Pizza. Poder sumar y restar vectores y obtener resultados con sentido refleja una estructura lineal.

\textit{«Esto es, en esencia, lo fundamental que está ocurriendo: la dirección importa.»} Todo lo que se ha observado hasta ahora en las redes neuronales naturales es consistente con la hipótesis de la representación lineal.
\end{importantbox}

Chris también mencionó trabajo reciente sobre representaciones no lineales (características multidimensionales / variedades), pero su actitud hacia la hipótesis lineal es pragmática: citó una analogía espléndida:

\begin{knowledgebox}{La lección de la teoría del calórico}
Incluso una teoría equivocada (como la teoría del calórico), si se toma en serio y se lleva al extremo, puede producir resultados prácticos: el motor de combustión fue desarrollado por personas que creían en la teoría del calórico.

\textit{«Vale la pena tomar en serio una hipótesis y llevarla al límite.»} Aunque la hipótesis de la representación lineal termine por no ser del todo correcta, avanzar dentro de su rango de validez puede producir descubrimientos importantes.
\end{knowledgebox}

\subsection{La hipótesis de Superposition}

Un embedding de palabras de 500 dimensiones no puede corresponder solo a 500 conceptos: el modelo necesita representar muchos más conceptos que dimensiones tiene. La solución: aprovechar las propiedades geométricas del espacio de alta dimensión más la dispersión de los conceptos («Japón» e «Italia» rara vez aparecen juntos).

\begin{importantbox}{«La sombra de una red más grande y dispersa»}
\textit{«Las redes neuronales podrían ser la sombra de redes neuronales más grandes y dispersas; lo que vemos son esas proyecciones.»}

Imagina un «modelo de arriba» —enorme pero disperso, con cada neurona interpretable—. La red neuronal real es una proyección comprimida de ese modelo de arriba. Aprender = construir una compresión eficiente del modelo de arriba. El descenso de gradiente podría estar buscando en secreto en el espacio de modelos extremadamente dispersos, para después plegarlos en una matriz densa.

Cantidad de conceptos: el lema de Johnson-Lindenstrauss muestra que el número de direcciones aproximadamente ortogonales que se pueden incorporar es exponencial respecto al número de dimensiones. La dispersión y la estructura de correlación aumentan aún más esta cifra.
\end{importantbox}

\subsection{El avance de Sparse Autoencoder}

Si la hipótesis de Superposition es correcta, Dictionary Learning / Sparse Autoencoder es la solución natural.

\textbf{Octubre de 2023}, «Towards Monosemanticity»: al usar Sparse Autoencoder en un modelo de una sola capa, \textit{«características bellas e interpretables surgieron así de manera natural»}. Se descubrieron: características del árabe, características del hebreo, características de Base64, características de lenguajes de programación, un «the» específico de contexto (que aparece en contextos matemáticos y predice «vector» y «matrix»), y patrones alternantes de medios caracteres Unicode.

Al entrenar el modelo dos veces, se encontraron características similares en ambos; Universality quedó confirmada de nuevo. \textit{«Una técnica muy natural simplemente funcionó... esta es en realidad una situación muy buena.»}

\textbf{Mayo de 2024}: se extendió a Claude 3 Sonnet (un modelo de producción). Se requería una \textbf{cantidad enorme} de GPU; las Scaling Laws for Interpretability de Tom Henighan ayudaron a predecir el tamaño óptimo del Sparse Autoencoder y el número óptimo de tokens de entrenamiento.

\begin{warningbox}{Características multimodales abstractas «escalofriantes»}
En Claude 3 Sonnet se descubrieron características fascinantes pero también inquietantes:

\begin{itemize}
    \item \textbf{Característica de vulnerabilidad de seguridad}: responde tanto al texto «disable SSL» \textbf{como} a imágenes de alguien haciendo clic en una advertencia de SSL de Chrome
    \item \textbf{Característica de puerta trasera}: responde tanto a puertas traseras en código \textbf{como} a imágenes de dispositivos de cámara oculta
    \item \textbf{Característica de engaño/mentira}: al forzar su activación, Claude comienza a mentir
    \item Características de búsqueda de poder, de dar un golpe de Estado, de ocultar información, entre otras
\end{itemize}

\textit{«Esto muestra cuán abstractos son estos conceptos.»} La existencia de estas características se relaciona directamente con los requisitos de seguridad de ASL-4 que Dario discutió en la Parte I: se necesita Interpretability para detectar si el modelo está engañando.
\end{warningbox}

\subsection{Los límites de la interpretabilidad automatizada}

Una idea natural es usar Claude para etiquetar las características descubiertas por Sparse Autoencoder. Chris reconoce que la interpretabilidad automatizada tiene valor, pero \textit{«tiene sus dudas al respecto»}:

Las etiquetas que da la IA son \textit{«en cierto sentido verdaderas, pero no capturan realmente la característica específica»}. Esto es similar al escepticismo de los matemáticos ante las demostraciones automatizadas por computadora: crees que la conclusión es correcta, pero no \textbf{comprendes}.

Una preocupación de seguridad más profunda: el clásico artículo de Ken Thompson «Reflections on Trusting Trust»: si usas IA para verificar la seguridad de otra IA, ¿puedes confiar en el auditor? Por ahora esto no es un gran problema, pero a medida que los sistemas se vuelvan más poderosos, esta pregunta de «quién vigila a los vigilantes» se volverá crucial.

\subsection{Direcciones futuras: de lo microscópico a lo macroscópico}

\textbf{De las características a los circuitos}: hasta ahora se ha entendido la representación (qué se activa), pero no el proceso de cómputo (cómo fluye y se transforma la información). Los pesos de interferencia (un artefacto de Superposition) hacen más difícil el análisis de circuitos.

El problema de la \textbf{«materia oscura»}: los Sparse Autoencoder actuales solo alcanzan a ver una pequeña parte de la «materia» de la red neuronal. \textit{«Es como la astronomía temprana: a medida que construimos mejores Sparse Autoencoder... vemos cada vez más estrellas.»}

\textbf{De la microbiología a la anatomía}: el Mech Interp actual es «la microbiología de las redes neuronales», algo extremadamente fino. Pero las preguntas que nos interesan son macroscópicas. Es necesario ascender por los niveles: biología molecular $\to$ biología celular $\to$ histología $\to$ anatomía $\to$ zoología $\to$ ecología. O, en la analogía de la física: partículas individuales $\to$ física estadística $\to$ termodinámica.

\begin{knowledgebox}{¿Las redes neuronales tienen «órganos»?}
\textit{«Espero que exista algo mucho más grande que las características y los circuitos.»} ¿Existen estructuras macroscópicas similares al corazón, a regiones cerebrales o al sistema respiratorio?

No se puede saltar directamente a lo macroscópico: primero hay que entender la estructura microscópica y después estudiar los patrones de conexión. Es como no poder saltarse la biología molecular para estudiar directamente la anatomía.
\end{knowledgebox}

\subsection{Redes neuronales artificiales frente a biológicas}

Chris considera que el trabajo de los neurocientíficos es \textit{«mucho más difícil»}. La lista de ventajas de estudiar redes neuronales artificiales es impresionante:

\begin{itemize}
    \item Registrar \textbf{todas} las neuronas (no solo las accesibles)
    \item Introducir cualquier dato de entrada
    \item Las neuronas no cambian durante el estudio
    \item Se puede hacer ablación (eliminar) de cualquier neurona
    \item Se puede editar cualquier peso de conexión
    \item Se puede deshacer cualquier cambio
    \item Se puede forzar la activación de cualquier neurona a cualquier valor
    \item Se conoce el conectoma completo, \textbf{incluidos los pesos exactos} (no solo conexiones binarias)
    \item Se pueden calcular gradientes
\end{itemize}

\textit{«Tenemos tantas ventajas sobre los neurocientíficos, y aun con todas esas ventajas, esto sigue siendo muy difícil. Si para nosotros es tan difícil, bajo las restricciones de la neurociencia parece casi imposible.»}

Esta es también la razón por la que Chris contrata activamente personas del campo de la neurociencia: \textit{«un problema más fácil, pero aun así muy difícil»}.

\subsection{Seguridad y belleza: el doble llamado de Mech Interp}

Algunas personas se sienten \textit{«decepcionadas»} con las redes neuronales: \textit{«solo son reglas simples amplificadas.»} La réplica de Chris es sumamente elegante: la evolución también es \textit{«solo reglas simples»}: mutación aleatoria más selección natural. Y sin embargo produjo todo el esplendor de la biología.

\textit{«La belleza está en que la simplicidad produce complejidad.»} En el interior de las redes neuronales se ha creado \textit{«una enorme complejidad y belleza que la gente normalmente no ve. Una estructura increíblemente rica que espera ser descubierta.»}

La imagen final: los circuitos crecen hacia la luz de la función de pérdida. \textit{«Este organismo que hemos cultivado, y no sabemos qué es lo que hemos cultivado.»} Esto es, a la vez, el motor científico de Mech Interp y su motor estético.

Para cerrar, en palabras de Alan Watts: \textit{«The only way to make sense out of change is to plunge into it, move with it, and join the dance.»}

\subsection{Resumen de la sección}

La hipótesis de la representación lineal y la hipótesis de Superposition constituyen la base teórica actual de Mech Interp. Sparse Autoencoder extrae características interpretables a partir de Superposition; la extensión exitosa desde modelos de una sola capa hasta Claude 3 Sonnet demuestra la viabilidad del método. La interpretabilidad automatizada es útil, pero no puede sustituir la comprensión humana. El reto central para el futuro es pasar de la «microbiología» a la «anatomía»: descubrir la estructura macroscópica de las redes neuronales. Mech Interp no solo sirve a la seguridad, sino que también explora la profunda pregunta científica de «qué es exactamente lo que hemos creado».
\section{La hoja de ruta de Anthropic: cómo convergen las tres líneas}

\subsection{De Scaling al producto, y de ahí a la gobernanza}

Lo más valioso de esta entrevista tan extensa no es solo lo que cada uno de los tres invitados expresó por separado, sino que muestra cómo, dentro de un laboratorio de vanguardia, se entretejen en una sola hoja de ruta las tres líneas de \textbf{capacidad, comportamiento y gobernanza}.

\begin{table}[H]
\centering
\begin{tabular}{p{0.18\textwidth} p{0.23\textwidth} p{0.25\textwidth} p{0.22\textwidth}}
\toprule
Nivel de la hoja de ruta & Persona representativa & Pregunta central & Mecanismo correspondiente de Anthropic \\
\midrule
Nivel de capacidad & Dario & Si el modelo sigue volviéndose más fuerte y cuándo cruzará umbrales clave & Scaling, posentrenamiento, Computer Use, la cronología de una IA poderosa \\
Nivel de comportamiento & Amanda & Cómo se comporta el modelo de manera confiable, estable y sin adulación excesiva en las interacciones cotidianas & la personalidad de Claude, CAI, el System Prompt, evaluaciones de comportamiento \\
Nivel verificable & Chris & Cómo sabemos qué está haciendo realmente el modelo por dentro & SAE, Monosemanticity, interpretabilidad automatizada, detección de engaño \\
Nivel de gobernanza & Acción conjunta de los tres & Cómo limitar el riesgo cuando la capacidad se acerca a umbrales peligrosos & RSP, ASL, evaluaciones externas, iniciativas de política \\
\bottomrule
\end{tabular}
\caption{La hoja de ruta implícita de cuatro niveles de Anthropic en la entrevista}
\end{table}

\begin{importantbox}{No son tres intereses distintos, sino un ciclo cerrado}
Dario se encarga de llevar la capacidad a la vanguardia, Amanda se encarga de moldear al modelo como un agente conductual utilizable por la sociedad, y Chris se encarga de que ese sistema sea verificable. Ninguna de estas líneas es un adorno prescindible; juntas constituyen la narrativa organizacional que distingue con más claridad a Anthropic frente a otros laboratorios de vanguardia.
\end{importantbox}

\subsection{Por qué el diseño de la personalidad y la interpretabilidad no son dos cosas distintas}

En apariencia, Amanda habla del tono, la adulación y la personalidad de Claude, mientras que Chris habla de SAE y de la descomposición en características, como si fueran mundos completamente distintos; pero la entrevista en realidad sugiere que terminarán por converger.

\begin{itemize}
    \item Si solo podemos observar el comportamiento del modelo a través de sus salidas, entonces el diseño de la personalidad solo puede quedarse en el nivel de «parches de comportamiento explícito»
    \item Si la interpretabilidad logra identificar de manera estable ciertas tendencias internas, como el engaño, la adulación o la búsqueda de poder, entonces la personalidad y la seguridad podrán entrar en una etapa verificable
    \item Dario menciona que ASL-4 quizá necesite usar la interpretabilidad para detectar sandbagging, y este es precisamente un ejemplo típico de la convergencia de las dos líneas
\end{itemize}

\begin{knowledgebox}{La personalidad de Claude puede verse como la «capa visible de la alineación»}
El System Prompt, el estilo para rechazar solicitudes, el tono de apoyo, la conducta no adulatoria: estos son resultados visibles para el usuario; mientras que el feature/circuit que le interesa a Chris podría ser el mecanismo causal interno detrás de esos resultados. Dicho de otro modo, Amanda moldea la interfaz y Chris busca las variables subyacentes.
\end{knowledgebox}

\subsection{Los puntos de inflexión clave entre 2025 y 2027}

Si se acepta el juicio de Dario, entonces lo más importante en los próximos dos o tres años no es un benchmark aislado, sino si los siguientes puntos de inflexión ocurren de manera simultánea:

\begin{enumerate}
    \item \textbf{Que el modelo supere el nivel de un experto humano en tareas laborales reales}: sobre todo en programación, biología y asistencia a la investigación
    \item \textbf{Que la capacidad de los Agentes aumente de manera notable}: pasar de responder preguntas en turnos cortos a ejecutar tareas que duran horas o incluso días
    \item \textbf{Que los umbrales de seguridad empiecen a activarse de manera institucionalizada}: que ASL-3 o un nivel superior deje de ser una política abstracta y se convierta en un umbral real para el despliegue
    \item \textbf{Que la interpretabilidad pase de ser un proyecto de investigación a un requisito operativo}: que deje de ser solo artículos y se vuelva parte del proceso de lanzamiento seguro
\end{enumerate}

\begin{warningbox}{El modo de falla más peligroso no es que «el modelo pierda el control de repente», sino un error de juicio organizacional}
En la entrevista aparece repetidamente un riesgo más realista: que, mientras la capacidad del modelo da saltos rápidos, sigamos gestionándolo con supuestos organizacionales antiguos, procesos de revisión antiguos y una mentalidad de producto antigua. Así, incluso si el modelo mismo no llega a perder el control al estilo de la ciencia ficción, el rezago institucional amplificaría el riesgo de todas formas.
\end{warningbox}

\subsection{Implicaciones para investigadores y constructores}

Esta entrevista no ofrece una conclusión única, sino una agenda de investigación:

\begin{itemize}
    \item Si te interesa la frontera de la capacidad, debes estudiar scaling, post-training, tool use y las aplicaciones de alto apalancamiento
    \item Si te interesa la posibilidad de desplegar el modelo, debes estudiar sycophancy, la consistencia de la personalidad, las evaluaciones y la relación humano-máquina
    \item Si te interesa la seguridad a largo plazo, debes adentrarte en mechanistic interpretability, la detección de engaño y la verificación independiente
    \item Si te interesan las consecuencias sociales, no puedes pasar por alto la regulación, la inercia institucional y los problemas de distribución del poder
\end{itemize}

\subsection{Resumen de la sección}

Lo distintivo de Anthropic no es que trabaje simultáneamente en el modelo, el producto y la seguridad, sino que intenta convertir las tres cosas en una sola ingeniería de sistemas. Aquí la capacidad, la personalidad, la interpretabilidad y la gobernanza no son proyectos paralelos, sino un ciclo cerrado en el que se limitan y se refuerzan mutuamente.

\section{Síntesis y ampliación}

\subsection{Ideas centrales de los tres invitados}

\begin{table}[H]
\centering
\begin{tabular}{lp{10cm}}
\toprule
\textbf{Invitado} & \textbf{Idea central} \\
\midrule
Dario Amodei & Scaling continuará; una IA poderosa podría llegar entre 2026 y 2027; el mayor riesgo es la concentración de poder \\
Amanda Askell & La personalidad de Claude es un proyecto de alineación; el marco del «viajero del mundo»; un buen carácter importa más que las reglas \\
Chris Olah & Las redes neuronales son «organismos que crecen»; Superposition es el enigma central; los SAE son la herramienta que produjo el avance \\
\bottomrule
\end{tabular}
\end{table}

\subsection{Relaciones entre temas}

Los tres forman un panorama completo: la RSP de Dario depende de que la interpretabilidad de Chris detecte el engaño (ASL-4); el diseño de personalidad de Amanda interactúa con la CAI y el RLHF; el «rasgo de engaño» de Chris se relaciona directamente con el riesgo de autonomía que le preocupa a Dario.

\subsection{Los cinco juicios de la entrevista que más vale meditar}

\begin{enumerate}
    \item \textbf{Scaling sigue siendo el eje central, pero el verdadero cuello de botella podría desplazarse hacia la velocidad de adaptación institucional y organizacional}
    \item \textbf{El posentrenamiento no solo embellece las salidas, sino que libera capacidades ya existentes en forma de interfaces utilizables}
    \item \textbf{El diseño de personalidad no es un retoque de la interfaz de usuario, sino parte de la ingeniería de alineación}
    \item \textbf{Para que la interpretabilidad mecanicista sea verdaderamente relevante, debe pasar del flujo de trabajo de los artículos académicos al flujo de trabajo de lanzamiento y auditoría}
    \item \textbf{El riesgo más profundo no necesariamente proviene del modelo en sí, sino de cómo se concentra y se usa el poder}
\end{enumerate}

\subsection{Divergencias y tensiones implícitas entre los tres invitados}

Esta entrevista no consiste en que las tres personas repitan la misma idea. Lo más interesante es que sus focos de atención son distintos, y eso constituye precisamente la tensión que existe de verdad al interior de las organizaciones de IA de vanguardia.

\begin{table}[H]
\centering
\begin{tabular}{p{0.18\textwidth} p{0.24\textwidth} p{0.24\textwidth} p{0.22\textwidth}}
\toprule
Tema & Foco de Dario & Foco de Amanda & Foco de Chris \\
\midrule
Avance del modelo & Scaling continuo, posentrenamiento y conversión en agentes & Si el comportamiento es aceptable para la sociedad & Si realmente entendemos los mecanismos internos del modelo \\
Método de seguridad & RSP, ASL, umbrales institucionales & Personalidad, coherencia, conducta no aduladora & Interpretabilidad, detección de engaño, verificación a nivel de rasgos \\
Riesgo a largo plazo & Concentración de poder, riesgo de autonomía, rezago institucional & Relación humano-IA, honestidad y empatía & Opacidad de los mecanismos internos que genera falta de control \\
Trabajo más prometedor & Biología, programación, posentrenamiento, despliegue & Alineación, diseño de personalidad, evaluación & SAE, monosemanticity, interpretación automatizada \\
\bottomrule
\end{tabular}
\caption{Diferencias y complementariedades entre los focos de los tres invitados}
\end{table}

\begin{knowledgebox}{La verdadera dificultad no es quién tiene la razón, sino cómo poner en práctica las tres perspectivas al mismo tiempo}
Si una organización solo escucha a Dario, es fácil que sobrestime el avance de capacidades y subestime los detalles de comportamiento; si solo escucha a Amanda, es fácil que limite la alineación a la capa de interfaz; si solo escucha a Chris, puede quedarse durante mucho tiempo en las aguas profundas de la investigación sin cerrar el ciclo hacia el producto. El valor de la entrevista está precisamente en poner las tres perspectivas en la misma mesa.
\end{knowledgebox}

\subsection{Lista de acciones para lectores con distintos roles}

En realidad, esta entrevista deja enseñanzas distintas según el rol del lector:

\begin{itemize}
    \item \textbf{Investigadores de modelos}: prestar atención sobre todo a la ventana de oportunidad de Scaling, el posentrenamiento, la interpretabilidad y la AI for science
    \item \textbf{Ingenieros de producto}: prestar atención sobre todo a la sycophancy, el System Prompt, la coherencia de personalidad y el aspecto de ataque y defensa de Computer Use
    \item \textbf{Investigadores de seguridad y políticas públicas}: prestar atención sobre todo a la RSP, la ASL, el rezago institucional, la precisión regulatoria y la distribución del poder
    \item \textbf{Emprendedores o responsables de organizaciones}: prestar atención sobre todo a la densidad de talento, la destreza cultural, y cómo hacer que la seguridad sea una dimensión competitiva y no un centro de costos
\end{itemize}

\begin{importantbox}{Esta no es una entrevista que pertenezca solo a Anthropic}
Aunque los tres invitados vienen de Anthropic, el valor de este contenido no está en el chisme corporativo, sino en que revela cuatro preguntas que enfrentan en general las organizaciones de IA de vanguardia: cómo avanzar en capacidades, cómo moldear el comportamiento, cómo verificar internamente y cómo la sociedad puede absorberlo. Ya sea que trabajes en OpenAI, Google, Meta o la comunidad de código abierto, estas cuatro preguntas siempre regresan.
\end{importantbox}

\subsection{Implicaciones estratégicas para distintos tipos de organizaciones}

Si se ubica esta entrevista de nuevo en el mapa de la industria, se puede ver además cómo debería leerla cada tipo de organización:

\begin{table}[H]
\centering
\begin{tabular}{p{0.2\textwidth} p{0.28\textwidth} p{0.28\textwidth}}
\toprule
Tipo de organización & La parte que más vale asimilar & El riesgo que más se suele pasar por alto \\
\midrule
Empresas de modelos de vanguardia & Scaling, posentrenamiento, RSP, la hoja de ruta de interpretabilidad & Perseguir solo la capacidad e ignorar la coordinación entre comportamiento y gobernanza \\
Startups de la capa de aplicaciones & La personalidad de Claude, Computer Use, la conversión de flujos de trabajo en producto & Subestimar el impacto de saltos repentinos en la capacidad del modelo subyacente sobre los límites del producto \\
Equipos técnicos empresariales & La interpretabilidad, los límites de permisos, la revisión y el ritmo de despliegue & Confundir la mejora en la capacidad del modelo con una mejora en la readiness de la organización \\
Instituciones de políticas públicas y gobernanza & La ASL, el rezago institucional, el problema de la concentración de poder & Quedarse en discusiones abstractas sin condiciones de activación concretas ni herramientas de verificación \\
\bottomrule
\end{tabular}
\caption{Los puntos clave que cada tipo de organización debería extraer de esta entrevista}
\end{table}

Dicho de otro modo, esta entrevista no es solo una autoexplicación de Anthropic, sino también un caso de estudio sobre el diseño de organizaciones de IA de vanguardia: nos dice que la capacidad, el comportamiento, la verificación y la gobernanza deben modelarse al mismo tiempo, porque, de lo contrario, cualquier línea que avance sola puede convertirse en un punto débil sistémico.

\subsection{Rutas de lectura para seguir ampliando}

Si se toma esta entrevista como un mapa, la lectura posterior puede profundizar a lo largo de tres líneas:

\begin{itemize}
    \item \textbf{Línea de capacidad}: Scaling Laws, RLHF, Computer Use, AI for Biology
    \item \textbf{Línea de comportamiento}: Constitutional AI, sycophancy, model personality, human-AI relationship
    \item \textbf{Línea de verificación}: SAE, monosemanticity, sandbagging detection, AI control
\end{itemize}

\subsection{Cierre final}

Si hay que resumir en una sola frase esta entrevista de cinco horas, sería esta: \textbf{la verdadera competencia de la IA de vanguardia ya no consiste solo en qué modelo es más fuerte, sino en quién puede colocar un modelo más fuerte dentro de un sistema organizacional más estable.} Dario aporta la curva de capacidad, Amanda discute los límites de comportamiento aceptables para la sociedad y Chris intenta establecer la verificabilidad interna; juntos, los tres constituyen una hoja de ruta que tiene posibilidades de sostenerse a largo plazo.

\subsection{La conclusión más práctica para el lector}

Para el lector común, esta entrevista termina por reducirse a tres juicios muy prácticos:

\begin{enumerate}
    \item No entender la IA de vanguardia solo a través de las «tablas de clasificación de modelos»; hay que ver al mismo tiempo el comportamiento, la verificación y la gobernanza
    \item No entender la alineación como un tono cortés o una estrategia para negarse a responder; al final debe conectarse con los mecanismos internos y el diseño institucional
    \item No separar el riesgo a largo plazo de la conversión en producto en el mundo real; en la narrativa de Anthropic, ambos son en realidad dos caras de una misma cosa
\end{enumerate}

\subsection{Matriz de capacidad-comportamiento-verificación-gobernanza}

\begin{table}[H]
\centering
\begin{tabular}{p{0.18\textwidth} p{0.2\textwidth} p{0.2\textwidth} p{0.2\textwidth}}
\toprule
Dimensión & Pregunta representativa en la entrevista & Lo que exige de la organización & Qué pasa si falta \\
\midrule
Capacidad & Si el modelo todavía puede seguir con el scaling, si todavía puede completar tareas más largas & Un ritmo continuo de entrenamiento, posentrenamiento y conversión en producto & Quedar rápidamente rezagado frente a modelos más fuertes \\
Comportamiento & Si el modelo es honesto, no adulador y aceptable para el usuario & Diseño de personalidad, CAI, evaluación y estrategia de interacción & Un deterioro rápido de la usabilidad del producto y de la confianza \\
Verificación & Si sabemos qué está haciendo el modelo por dentro & Investigación de interpretabilidad, seguimiento de rasgos, verificación de seguridad independiente & Que el juicio de seguridad dependa en exceso del comportamiento externo \\
Gobernanza & Cuándo se activan las restricciones, quién es responsable, cómo evitar la concentración de poder & RSP, ASL, coordinación entre la organización y las políticas públicas & Que el crecimiento de la capacidad sea más rápido que la preparación institucional \\
\bottomrule
\end{tabular}
\caption{La matriz de cuatro dimensiones que resume esta entrevista}
\end{table}

El sentido de esta matriz es que comprime de verdad la discusión de los tres invitados en un marco ejecutable. Muchas organizaciones invierten solo en una de estas dimensiones, por ejemplo, invertir solo en la capacidad del modelo, hacer solo parches al comportamiento del producto, o limitarse a declaraciones de política, pero la entrevista sugiere una y otra vez: si falta cualquiera de estas dimensiones, el sistema se desequilibra.

\subsection{Lecturas adicionales}

\begin{itemize}
    \item Dario Amodei, \textit{Machines of Loving Grace}, 2024
    \item Amanda Askell, \textit{The optimal rate of failure}, Blog
    \item Chris Olah et al., \textit{Towards Monosemanticity}, Anthropic, 2023
    \item Chris Olah et al., \textit{Scaling Monosemanticity}, Anthropic, 2024
    \item Anthropic, \textit{Responsible Scaling Policy}, 2023
\end{itemize}

\end{document}
